\chapter{التوابع وحركة الكواكب}\label{ch:g12:satellites-kepler}

الطبق على الشرفة مثبَّت بالبراغي، ومع ذلك يحدّق طول اليوم في نقطة
واحدة ثابتة من السماء، على ارتفاع \qty{36000}{km}؛ ومحطة الفضاء
تعبر سماء المساء في دقائق؛ وكوكبة من الساعات الطائرة
تخبر هاتفك أين تقف. وليس في أي من هذه الآلات
محرك يعمل: فكلها ببساطة \emph{تسقط}. ويحسب هذا الفصل
كم يجب أن تسقط كلٌّ منها سرعةً وارتفاعًا --- ثم يزن
الكواكب بالقانون نفسه.

\section{السقوط حول الأرض}

كل شيء يقوم على القانون المقيس سابقًا
(\cref{ch:g10:universal-gravitation}): فالأرض، وكتلتها $M$، تجذب
تابعًا كتلته $m$ ويبعد مركزه عن مركزها $r$
\omterm{def:g10:inertia:force}{بقوة} موجَّهة نحو مركز الأرض، وشدتها
\[ F = G\,\frac{m M}{r^2}, \qquad G = \qty{6.67e-11}{N.m^2/kg^2}. \]
وانتبه إلى ما هو $r$: إنها المسافة \emph{بين المركزين}. \omterm{def:g10:universal-gravitation:satellite}{فالتابع} على
الارتفاع $h$ له $r = R + h$، حيث $R = \qty{6.37e6}{m}$ نصف قطر
الأرض --- ونسيان $R$ هو الخطأ الكلاسيكي في هذا الفصل.

\begin{remark}[مدفع نيوتن، من جديد]\label{rem:g12:satellites-kepler:cannon}
إذا أُطلقت كرة مدفع أفقيًّا هبطت؛ وإذا أُطلقت أسرع هبطت أبعد؛
وعند \omterm{def:g12:kinematics-2d:velocity}{السرعة} الصحيحة تنحني الأرض مبتعدة \omterm{def:g12:kinematics-2d:velocity}{بالسرعة} نفسها التي تسقط بها
الكرة، فينغلق السقوط على نفسه: أي \omterm{def:g10:universal-gravitation:satellite}{مدار}. \omterm{def:g10:universal-gravitation:satellite}{فالتابع}
لا يحتاج محركًا --- إنما يحتاج \omterm{def:g12:kinematics-2d:velocity}{السرعة} \emph{الجانبية} الصحيحة.
\end{remark}

\begin{omfigure}
\begin{tikzpicture}[scale=1.0, font=\small]
  \fill[omExo!15] (0,0) circle (1.5); \draw[omExo, thick] (0,0) circle (1.5);
  \fill[omExo!60] (-0.22,1.44) -- (0,1.85) -- (0.22,1.44) -- cycle;
  \coordinate (L) at (0,1.85);
  \draw[omProp, thick] (L) .. controls (0.7,1.85) .. (50:1.5);
  \draw[omProp, thick] (L) .. controls (1.4,1.85) and (2.0,0.9) .. (0:1.5);
  \draw[-Stealth, omDef, thick] (0,1.85) arc (90:-260:1.85);
  \draw[-Stealth, omThm, thick] (L) .. controls (1.7,2.0) and (3.2,2.2)
    .. (4.3,2.9);
  \node[omDef, anchor=east] at (-1.85,1.1) {مدار};
  \node[omThm, anchor=west] at (3.35,2.35) {إفلات};
  \node[omProp, anchor=west] at (1.35,0.55) {يسقط};
\end{tikzpicture}

{\small مدفع نيوتن: كلما أُطلقت الكرة أسرع هبطت أبعد
(بالبرتقالي)؛ وعند نحو \qty{7.9}{km/s} يوافق سقوطها تقوّس
الأرض فتدور في \omterm{def:g10:universal-gravitation:satellite}{مدار} (بالأزرق)؛ وأسرع من ذلك تفلت (بالأحمر).}
\end{omfigure}

\section{المدار الدائري: سرعة واحدة ودور واحد}

\begin{theorem}[السرعة المدارية]\label{thm:g12:satellites-kepler:speed}
\omterm{def:g10:universal-gravitation:satellite}{تابع} في \omterm{def:g10:universal-gravitation:satellite}{مدار} دائري منتظم نصف قطره $r$ حول جسم كتلته
$M$ ينتقل \omterm{def:g12:kinematics-2d:velocity}{بالسرعة}
\[ v = \sqrt{\frac{G M}{r}}, \]
وهي لا تتوقف على كتلة \omterm{def:g10:universal-gravitation:satellite}{التابع}.
\end{theorem}

\begin{proof}
\omterm{def:g11:fundamental-interactions:four}{الجاذبية} هي \omterm{def:g10:inertia:force}{القوة} الوحيدة، ومن ثم يُقرأ \omterm{thm:g12:newtons-laws:second}{قانون نيوتن الثاني}
(\cref{thm:g12:newtons-laws:second}) هكذا: $m\vect a = \vect F$؛
فأسقِطه على \omterm{def:g12:kinematics-2d:frenet}{معلم فريني} (\cref{def:g12:kinematics-2d:frenet}).
\omterm{def:g10:inertia:force}{والقوة} مصوَّبة نحو المركز: فمركّبتها المماسّية معدومة، ومن ثم
تكون \omterm{def:g12:kinematics-2d:velocity}{السرعة} ثابتة؛ ومركّبتها الناظمية هي $GmM/r^2$ كلها،
وفي الحركة الدائرية $a_n = v^2/r$
(\cref{thm:g12:kinematics-2d:centripetal}). ومنه $m v^2/r =
GmM/r^2$، وبقسمة $m$ --- $v^2 = GM/r$.
\end{proof}

\begin{omfigure}
\begin{tikzpicture}[scale=1.0, font=\small]
  \fill[omExo!15] (0,0) circle (0.85); \draw[omExo, thick] (0,0) circle (0.85);
  \node at (0,0) {الأرض}; \draw[dashed, omExo] (0,0) circle (2.3);
  \draw[thin] (0,0) -- (40:2.3) node[midway, above left] {$r$};
  \fill (40:2.3) circle (2.2pt) node[above right] {$S$};
  \draw[-Stealth, omThm, very thick] (40:2.3) -- (40:1.35)
    node[pos=0.85, below right] {$\vect F$};
  \draw[-Stealth, omDef, thick] (40:2.3) -- ++(130:1.15)
    node[above] {$\vect v$};
\end{tikzpicture}

{\small آلية \omterm{def:g10:universal-gravitation:satellite}{المدار} كلها: \omterm{def:g12:kinematics-2d:velocity}{سرعة} مماسّة، \omterm{def:g10:inertia:force}{وقوة}
نحو المركز --- \omterm{def:g11:fundamental-interactions:four}{فالجاذبية} لا تبطّئ \omterm{def:g10:universal-gravitation:satellite}{التابع} قط، إنما
تحني مساره أبدًا.}
\end{omfigure}

\begin{proposition}[الدور المداري]\label{prop:g12:satellites-kepler:period}
يكمل \omterm{def:g10:universal-gravitation:satellite}{التابع} \omterm{def:g10:universal-gravitation:satellite}{مداره} في \omterm{def:g12:kinematics-2d:ucm}{الدور}
\[ T = \frac{2\pi r}{v} = 2\pi \sqrt{\frac{r^3}{G M}}. \]
\end{proposition}

\begin{proof}
الدورة الواحدة هي المحيط $2\pi r$ \omterm{def:g12:kinematics-2d:velocity}{بسرعة} ثابتة؛ فأدخِل
$v = \sqrt{GM/r}$: $T = 2\pi r\sqrt{r/(GM)} = 2\pi\sqrt{r^3/(GM)}$.
\end{proof}

ولا يحوي $v$ ولا $T$ المقدارَ $m$: فبرغي يهتزّ حتى ينفكّ من محطة
الفضاء يدور في \omterm{def:g10:universal-gravitation:satellite}{مدار} إلى جانبها. وكلاهما يهبط مع $r$:
\emph{فكلما ارتفع \omterm{def:g10:universal-gravitation:satellite}{المدار} بطؤ \omterm{def:g10:universal-gravitation:satellite}{التابع}}، سرعةً
($v \propto 1/\sqrt r$) وزاويةً ($T \propto r^{3/2}$).

\begin{example}[محطة الفضاء]\label{ex:g12:satellites-kepler:iss}
تحلّق محطة الفضاء الدولية على ارتفاع
$h \approx \qty{400}{km}$، ومنه $r = \num{6.37e6} + \num{4.0e5} =
\qty{6.77e6}{m}$؛ وبأخذ $GM = \num{6.67e-11} \times \num{5.97e24} =
\qty{3.98e14}{m^3/s^2}$ يكون $v = \sqrt{\num{3.98e14}/\num{6.77e6}}
\approx \qty{7.7e3}{m/s}$ و$T = 2\pi r/v \approx \qty{5.5e3}{s}
\approx \qty{92}{min}$: أي دورة حول الكوكب كل ساعة ونصف،
ونحو ستة عشر شروقًا في اليوم لروّاد الفضاء.
\end{example}

\begin{example}[القمر يسقط كالتفاحة]\label{ex:g12:satellites-kepler:moon}
القمر، على $r = \qty{3.84e8}{m}$، يدور في $T = \qty{27.3}{d} =
\qty{2.36e6}{s}$، ومنه $v = 2\pi r/T \approx \qty{1.02e3}{m/s}$ و
$a = v^2/r \approx \qty{2.7e-3}{m/s^2}$ --- وهو بالضبط $g/60^2$، و
يقع القمر على بعد $60$ نصف قطر أرضي: فالجذب الذي يُسقط التفاحة
يحني القمر، مضعَّفًا بمربع المسافة في المقام. وفي كل \omterm{def:g10:orders-of-magnitude:unit}{ثانية}
``يسقط'' القمر نحو \qty{1.4}{mm} نحونا، ويحمله كيلومتره
الجانبي حولنا --- وهو التحقّق الذي أجراه نيوتن سنة 1666.
\end{example}

\section{قوانين كبلر}

لم يخمّن نيوتن مربع المسافة في المقام: بل استخرجه من ثلاث
انتظامات استخلصها يوهانس كبلر، بين 1609 و1619،
من عشرين سنة من مواضع الكواكب التي رصدها تيخو براهي بالعين المجردة ---
وهي أفضل معطيات عالم ما قبل المقراب.

\begin{definition}[الإهليلج]\label{def:g12:satellites-kepler:ellipse}
\emph{الإهليلج}\index{إهليلج} مجموعة النقاط التي يكون مجموع
مسافاتها إلى نقطتين ثابتتين، هما \emph{البؤرتان}\index{بؤرة (إهليلج)}،
ثابتًا؛ ونصف أطول أقطاره هو
\emph{نصف المحور الأكبر}\index{نصف المحور الأكبر} $a$، والدائرة
هي الحالة التي تنطبق فيها البؤرتان. وأقرب نقاط \omterm{def:g10:universal-gravitation:satellite}{مدار} كوكب إلى الشمس وأبعدها
هما \emph{الحضيض}\index{حضيض شمسي}
و\emph{الأوج}\index{أوج شمسي}.
\end{definition}

\begin{theorem}[قوانين كبلر]\label{thm:g12:satellites-kepler:kepler}
\index{قوانين كبلر}%
في الكواكب الدائرة حول الشمس (وكتلتها $M$):
\begin{enumerate}
  \item كل \omterm{def:g10:universal-gravitation:satellite}{مدار} \omterm{def:g12:satellites-kepler:ellipse}{إهليلج} والشمس في إحدى بؤرتيه؛
  \item والقطعة الواصلة بين الشمس والكوكب تكنس مساحات متساوية في أزمنة متساوية؛
  \item ونسبة مربع \omterm{def:g12:kinematics-2d:ucm}{الدور} إلى مكعب \omterm{def:g12:satellites-kepler:ellipse}{نصف المحور الأكبر} هي
        نفسها لكل الكواكب: $T^2/a^3 = 4\pi^2/(G M)$.
\end{enumerate}
والقوانين نفسها تحكم أي عائلة من التوابع حول أي جسم
مركزي، حيث $M$ الكتلة المركزية. \admitted
\end{theorem}

\begin{remark}[ما نقدر أن نبرهن عليه هذه السنة]\label{rem:g12:satellites-kepler:proof}
في \omterm{def:g10:universal-gravitation:satellite}{مدار} \emph{دائري} يكون القانون الثالث لنا: فتربيع
\cref{prop:g12:satellites-kepler:period} يعطي
$T^2/r^3 = 4\pi^2/(GM)$ --- أي ثابتًا واحدًا لكل \omterm{def:g10:universal-gravitation:satellite}{تابع} للمركز
نفسه. أما كون الإهليلجات والمساحات المتساوية والصيغة العامة $a$ تنتج
عن مربع المسافة في المقام فيُبيَّن، بمزيد من التحليل، في مجلد السنة الأولى.
وفي الأثناء يمكن قراءة القانون الثاني: فالكوكب أسرع ما يكون عند \omterm{def:g12:satellites-kepler:ellipse}{الحضيض}،
حيث يجب أن تسرع قطعة الكنس القصيرة.
\end{remark}

\begin{omfigure}
\begin{tikzpicture}[scale=0.85, font=\small,
    declare function={re(\t)=2.025/(1+0.5*cos(\t));}]
  \draw[thick] plot[domain=0:360, samples=120, smooth]
    ({1.35+re(\x)*cos(\x)}, {re(\x)*sin(\x)});
  \fill[omDef!30] (1.35,0) -- plot[domain=-40:40, samples=25]
    ({1.35+re(\x)*cos(\x)}, {re(\x)*sin(\x)}) -- cycle;
  \fill[omDef!30] (1.35,0) -- plot[domain=175:185, samples=10]
    ({1.35+re(\x)*cos(\x)}, {re(\x)*sin(\x)}) -- cycle;
  \fill[omThm] (1.35,0) circle (2.6pt) node[below=2pt] {الشمس};
  \draw[-Stealth, omDef, thick]
    ({1.35+re(40)*cos(40)}, {re(40)*sin(40)}) -- ++(-0.55,0.42);
  \draw[-Stealth, omDef, thick]
    ({1.35+re(175)*cos(175)}, {re(175)*sin(175)}) -- ++(-0.12,-0.5);
  \node[right] at (2.75,0) {$P$}; \node[left] at (-2.72,0) {$A$};
\end{tikzpicture}

{\small قانون كبلر الثاني: في شهر واحد قرب \omterm{def:g12:satellites-kepler:ellipse}{الحضيض} $P$ يكنس
الكوكب قطاعًا قصيرًا سمينًا؛ وقرب \omterm{def:g12:satellites-kepler:ellipse}{الأوج} $A$، قطاعًا طويلًا نحيلًا
--- بمساحتين متساويتين: فهو يسرع إذا قرب، ويتباطأ إذا بعد.}
\end{omfigure}

\begin{omfigure}
\begin{tikzpicture}
\begin{axis}[omaxis, width=8.8cm, height=5.8cm,
    xmode=log, ymode=log,
    xlabel={نصف المحور الأكبر $a$ (وحدات فلكية)},
    ylabel={الدور $T$ (بالسنوات)},
    xmin=0.25, xmax=50, ymin=0.15, ymax=400]
  \addplot[omThm, thick] coordinates {(0.3,0.1643) (40,253)};
  \addplot[only marks, mark=*, mark size=1.7pt, omDef] coordinates
    {(0.387,0.241) (0.723,0.615) (1,1) (1.524,1.881) (5.203,11.86)
     (9.537,29.46) (19.19,84.02) (30.07,164.8)};
  \node[below right, font=\scriptsize] at (axis cs:0.387,0.241) {عطارد};
  \node[below right, font=\scriptsize] at (axis cs:1,1) {الأرض};
  \node[below right, font=\scriptsize] at (axis cs:5.203,11.86) {المشتري};
  \node[below right, font=\scriptsize] at (axis cs:30.07,164.8) {نبتون};
\end{axis}
\end{tikzpicture}

{\small الكواكب الثمانية على محاور لوغاريتمية مزدوجة: مستقيم ميله
$3/2$، أي\ $T^2 \propto a^3$ --- وهو قانون كبلر الثالث، الذي يُقرأ في
هذه الوحدات (الوحدات الفلكية والسنوات) ببساطة $T^2 = a^3$.}
\end{omfigure}

\section{المدارات في العمل}

\begin{definition}[المدار الثابت بالنسبة إلى الأرض]\label{def:g12:satellites-kepler:geo}
يكون \omterm{def:g10:universal-gravitation:satellite}{التابع} \emph{ثابتًا بالنسبة إلى الأرض}\index{مدار ثابت بالنسبة إلى الأرض} إذا
تعلّق فوق نقطة أرضية واحدة ثابتة. ويجب أن يكون \omterm{def:g10:universal-gravitation:satellite}{مداره}
دائريًّا استوائيًّا، ودوره \emph{يوم نجمي}\index{يوم نجمي}
واحد، $T = \qty{86164}{s}$ (\qty{23}{h}
\qty{56}{min} \qty{4}{s}): أي الزمن الذي تدور فيه الأرض دورة
\emph{بالنسبة إلى النجوم}. (أما اليوم \qty{24}{h} فهو بالنسبة إلى
الشمس، التي تتقدم الأرض نحوها درجةً في اليوم أيضًا.)
\end{definition}

\begin{example}[ارتفاع المدار الثابت بالنسبة إلى الأرض]\label{ex:g12:satellites-kepler:geoalt}
قانون كبلر الثالث، مقروءًا معكوسًا، يملي نصف القطر:
$r = \left(GMT^2/4\pi^2\right)^{1/3} =
\left(\num{3.98e14} \times (\num{86164})^2/4\pi^2\right)^{1/3}
\approx \qty{4.22e7}{m}$ --- أي ارتفاع $h = r - R \approx
\qty{3.58e7}{m} \approx \qty{35800}{km}$، يُقطع \omterm{def:g12:kinematics-2d:velocity}{بسرعة} $v = 2\pi r/T
\approx \qty{3.1}{km/s}$. فكل مرحّل ``يتعلّق ساكنًا'' يقع على
هذه الدائرة الواحدة فوق خط الاستواء --- ولا دائرة سواها.
\end{example}

\begin{example}[كوكبة الملاحة]\label{ex:g12:satellites-kepler:gps}
تحلّق منظومات الملاحة بالتوابع بنحو ثلاثين تابعًا على
$r \approx \qty{2.66e7}{m}$ (بارتفاع $\approx \qty{20200}{km}$)،
حيث يعطي \cref{prop:g12:satellites-kepler:period} \omterm{def:g12:kinematics-2d:ucm}{الدور} $T \approx
\qty{4.32e4}{s}$: أي نصف \omterm{def:g12:satellites-kepler:geo}{يوم نجمي}، ومن ثم يعيد كل \omterm{def:g10:universal-gravitation:satellite}{تابع} اقتفاء
أثره الأرضي يوميًّا. ويقارن مستقبِلك أزمنة وصول
إشارات ساعاتها بمدارات معلومة إلى \omterm{def:g10:orders-of-magnitude:unit}{المتر}.
\end{example}

\begin{method}[وزن عالم من تابعه]\label{met:g12:satellites-kepler:weigh}
\index{وزن كوكب}%
لقياس كتلة كوكب أو نجم: جِد أي \omterm{def:g10:universal-gravitation:satellite}{تابع} له
(قمرًا أو مسبارًا أو كوكبًا)، وقِس نصف قطر \omterm{def:g10:universal-gravitation:satellite}{مداره} $r$ ودوره
$T$، واقلب القانون الثالث:
\[ M = \frac{4\pi^2 r^3}{G\,T^2}. \]
وهذا لا يزن إلا الجسم \emph{المركزي} --- إذ تلاشت كتلة \omterm{def:g10:universal-gravitation:satellite}{التابع}
في \cref{thm:g12:satellites-kepler:speed}. وكل كتلة في بطاقة
المعطيات جاءت بهذه الطريقة: فالأرض من القمر، والشمس من
\omterm{def:g10:universal-gravitation:satellite}{مدار} الأرض نفسه.
\end{method}

\begin{remark}[الأعلى أبطأ --- لكنه أغلى]\label{rem:g12:satellites-kepler:energy}
بما أن $v = \sqrt{GM/r}$، يزحف القمر \omterm{def:g12:kinematics-2d:velocity}{بسرعة} \qty{1}{km/s} بينما تنطلق
المحطة \omterm{def:g12:kinematics-2d:velocity}{بسرعة} \qty{7.7}{km/s}. ومع ذلك يكون \omterm{def:g10:universal-gravitation:satellite}{المدار} الأعلى
\emph{أغلى}: فالصعود ضد \omterm{def:g11:fundamental-interactions:four}{الجاذبية} يخزّن \omterm{def:g10:energy-conservation:energy}{طاقة} كامنة
(وتُسوَّى الحسابات في \cref{ch:g12:work-and-energy})، و
\omterm{def:g10:signals-and-waves:oscilloscope}{الكسب} في $E_p$ يفوق الخسارة في $E_k$. ومن هنا يحرق صاروخ وقودًا
\emph{إلى الأمام} ليبلغ مدارًا أعلى أبطأ --- وتابعٌ يكبحه هواء رقيق
يتحلزن \emph{نازلًا} ويتسارع.
\end{remark}

\begin{remark}[بطاقة المعطيات]\label{rem:g12:satellites-kepler:data}
ما لم يُذكر خلاف ذلك: $G = \qty{6.67e-11}{N.m^2/kg^2}$؛ والأرض:
$M = \qty{5.97e24}{kg}$، $R = \qty{6.37e6}{m}$، \omterm{def:g12:satellites-kepler:geo}{واليوم النجمي}
\qty{86164}{s}؛ \omterm{def:g10:universal-gravitation:satellite}{ومدار} القمر: $r = \qty{3.84e8}{m}$، $T =
\qty{27.3}{d}$؛ والشمس: $M = \qty{1.99e30}{kg}$؛ والمسافة بين الأرض والشمس
\qty{1.496e11}{m}؛ والسنة الواحدة $= \qty{3.156e7}{s}$؛ والمريخ: $M =
\qty{6.42e23}{kg}$، $R = \qty{3.39e6}{m}$.
\end{remark}

\section{تمارين}

\begin{exercise}[$\star$]\label{exo:g12:satellites-kepler:1}
\omterm{def:g10:universal-gravitation:satellite}{تابع} كتلته \qty{1000}{kg} يقف على منصة الإطلاق، ثم يحلّق على
ارتفاع \qty{400}{km}. احسب جذب الأرض في الموضعين:
فأيّ نسبة مئوية تبقى في الأعلى --- ولماذا يطفو من على متنه؟
\end{exercise}

\begin{exercise}[$\star$]\label{exo:g12:satellites-kepler:2}
في \omterm{def:g10:universal-gravitation:satellite}{تابع} (نظري) يلامس سطح الأرض
($r = R$)، احسب \omterm{def:g12:kinematics-2d:velocity}{السرعة} المدارية \omterm{def:g12:kinematics-2d:ucm}{والدور}. وقارن ذلك بكرة
مدفع نيوتن (\cref{rem:g12:satellites-kepler:cannon}).
\end{exercise}

\begin{exercise}[$\star$]\label{exo:g12:satellites-kepler:3}
يُنقل \omterm{def:g10:universal-gravitation:satellite}{تابع} من \omterm{def:g10:universal-gravitation:satellite}{مدار} دائري نصف قطره $r$ إلى \omterm{def:g10:universal-gravitation:satellite}{مدار} نصف قطره
$2r$. فبأي عامل تتغير سرعته ودوره؟ وهل تتوقف
الأجوبة على كتلته؟
\end{exercise}

\begin{exercise}[$\star$]\label{exo:g12:satellites-kepler:4}
أعِد حساب أرقام المحطة: انطلاقًا من $r = \qty{6.77e6}{m}$، احسب
$v$ و$T$، وعدد المدارات المكتملة في \qty{24}{h}.
\end{exercise}

\begin{exercise}[$\star$]\label{exo:g12:satellites-kepler:5}
اذكر \omterm{thm:g12:satellites-kepler:kepler}{قوانين كبلر} الثلاثة. وأين يكون مذنّب أسرع ما يكون على إهليلجه
الشديد الاستطالة، وأيّ قانون يقول ذلك؟ وهل يتوقف \omterm{def:g10:signals-and-waves:period}{دور}
\omterm{def:g10:universal-gravitation:satellite}{تابع} على كتلته هو؟
\end{exercise}

\begin{exercise}[$\star\star$]\label{exo:g12:satellites-kepler:6}
زِن الأرض: انطلاقًا من نصف قطر \omterm{def:g10:universal-gravitation:satellite}{مدار} القمر ودوره (من بطاقة
المعطيات)، احسب كتلة الأرض وقارنها بقيمة البطاقة.
\end{exercise}

\begin{exercise}[$\star\star$]\label{exo:g12:satellites-kepler:7}
تقترح شركة تابعًا يتعلّق دائمًا فوق باريس
(على خط العرض $49^\circ$ شمالًا). انطلاقًا من جهة \omterm{def:g11:fundamental-interactions:four}{الجاذبية}، اشرح
لماذا يكون مركز كل \omterm{def:g10:universal-gravitation:satellite}{مدار} دائري مركز الأرض؛ واستنتج
أن الاقتراح مستحيل. وأين \emph{يقدر} \omterm{def:g10:universal-gravitation:satellite}{تابع} أن يتعلّق؟
\end{exercise}

\begin{exercise}[$\star\star$]\label{exo:g12:satellites-kepler:8}
يدور فوبوس حول المريخ على $r = \qty{9.38e6}{m}$ في $T = \qty{7}{h}
\qty{39}{min}$. استنتج كتلة المريخ؛ وقارنها ببطاقة المعطيات.
\end{exercise}

\begin{exercise}[$\star\star$]\label{exo:g12:satellites-kepler:9}
يدور المريخ حول الشمس على $a = \qty{2.279e11}{m}$ في \qty{687}{d}.
احسب $T^2/a^3$ للمريخ وللأرض (من بطاقة المعطيات)، وتحقّق من
قانون كبلر الثالث، واستنتج كتلة الشمس.
\end{exercise}

\begin{exercise}[$\star\star$]\label{exo:g12:satellites-kepler:10}
اختبار نيوتن للقمر: انطلاقًا من بطاقة المعطيات، احسب \omterm{def:g12:kinematics-2d:velocity}{سرعة} القمر و
تسارعه المركزي $a = v^2/r$. وبيّن أنه يساوي $g/60^2$،
وقل لماذا أقنع ذلك نيوتن بأن قانونًا واحدًا يحكم التفاحة والقمر.
\end{exercise}

\begin{exercise}[$\star\star$]\label{exo:g12:satellites-kepler:11}
بوحدات الشمس (الوحدات الفلكية والسنوات) يُقرأ قانون كبلر الثالث
$T^2 = a^3$: أي مستقيم المخطط اللوغاريتمي المزدوج في الدرس. ضع
عليه الكويكب سيريس، $a = \qty{2.77}{au}$، ومذنّب هالي،
$a = \qty{17.8}{au}$: واحسب الدورين. وقد مرّ هالي \omterm{def:g12:satellites-kepler:ellipse}{بالحضيض}
آخر مرة سنة 1986 --- فمتى يُنتظر عودته؟
\end{exercise}

\begin{exercise}[$\star\star\star$]\label{exo:g12:satellites-kepler:12}
يمسح هواء رقيق المحطة، ومع ذلك \emph{تزداد} سرعتها كلما
فقدت ارتفاعًا. احسب $v$ على الارتفاعين \qty{400}{km} و
\qty{350}{km}، ثم حُلّ المفارقة: \omterm{def:g10:inertia:inventory}{فالاحتكاك} يبطّئ الأشياء
--- فمن أين تأتي \omterm{def:g12:kinematics-2d:velocity}{السرعة} الزائدة؟
\end{exercise}

\begin{exercise}[$\star\star\star$]\label{exo:g12:satellites-kepler:13}
يجد مقراب كوكبًا خارجيًّا يدور حول نجمه على
$r = \qty{7.48e9}{m}$ كل \qty{3.50}{d}. زِن النجم،
\omterm{def:g10:orders-of-magnitude:unit}{بالكيلوغرام} وبوحدات الشمس. وأيّ كتلة --- كتلة النجم أم الكوكب ---
\emph{لا يمكن} الحصول عليها بهذه الطريقة، ولماذا؟
\end{exercise}

\begin{exercise}[$\star\star\star$]\label{exo:g12:satellites-kepler:14}
يريد مستوطنو المريخ مرحّلًا ``ثابتًا بالنسبة إلى المريخ'' فوق قاعدتهم.
ويدور المريخ بالنسبة إلى النجوم في \qty{24}{h} \qty{37}{min}:
احسب نصف قطر \omterm{def:g10:universal-gravitation:satellite}{المدار} والارتفاع فوق سطح
المريخ (من بطاقة المعطيات).
\end{exercise}

\begin{exercise}[$\star\star\star$]\label{exo:g12:satellites-kepler:15}
يمرّ مذنّب هالي \omterm{def:g12:satellites-kepler:ellipse}{بالحضيض} على $r_p = \qty{8.77e10}{m}$ \omterm{def:g12:kinematics-2d:velocity}{بسرعة}
\qty{54.5}{km/s}؛ وأوجه على $r_a = \qty{5.25e12}{m}$. وعند
الطرفين تكون \omterm{def:g12:kinematics-2d:velocity}{السرعة} عمودية على القطعة الواصلة بين الشمس والمذنّب،
ويفرض قانون كبلر الثاني عندئذٍ $r_p v_p = r_a v_a$
(بمثلثَي كنس متساويي المساحة $\tfrac12 r v\,\Delta t$). استنتج
\omterm{def:g12:kinematics-2d:velocity}{سرعة} \omterm{def:g12:satellites-kepler:ellipse}{الأوج} وعلّق على النسبة.
\end{exercise}

\section{مسألة: مقعد فوق خط الاستواء}

\begin{problem}[{مسألة نهاية الأسبوع --- وضع \omterm{def:g10:universal-gravitation:satellite}{تابع} فوق
خط الاستواء: لماذا يكون موقف السماء الوحيد دائرةً
على ارتفاع \qty{35800}{km}، وكيف يقارَن بالمحطة المنطلقة تحته،
وكيف يزن القانون نفسه المشتري في الطريق}]
\label{pb:g12:satellites-kepler:1}
تريد شركة اتصالات مرحّلًا لا تضطر أطباقها الأرضية، وقد ثُبّتت
مرة واحدة، إلى مطاردته أبدًا: أي تابعًا \omterm{def:g12:satellites-kepler:geo}{ثابتًا بالنسبة إلى الأرض}. وأنت
محلّل المهمة؛ فاستعمل بطاقة المعطيات. وللجزء IV: يدور قمر المشتري آيو
على $r = \qty{4.22e8}{m}$ في $T = \qty{1.77}{d}$، ويدور المشتري
بالنسبة إلى النجوم في \qty{9}{h} \qty{56}{min}.

\medskip
\noindent\textbf{الجزء I --- قواعد التعلّق ساكنًا.}
\begin{enumerate}
  \item يجب أن يبقى \omterm{def:g10:universal-gravitation:satellite}{التابع} فوق نقطة أرضية واحدة ثابتة. فماذا
        يجب أن يكون دوره --- ولماذا \omterm{def:g12:satellites-kepler:geo}{اليوم النجمي}
        \qty{86164}{s} لا \qty{24}{h}؟
  \item في الحركة الدائرية تتجه \omterm{def:g11:forces-and-motion:netforce}{القوة المحصّلة} نحو مركز
        الدائرة؛ وهنا \omterm{def:g10:inertia:force}{القوة} الوحيدة هي \omterm{def:g11:fundamental-interactions:four}{الجاذبية}. استنتج أن
        مركز أي \omterm{def:g10:universal-gravitation:satellite}{مدار} دائري هو مركز الأرض.
  \item ودائرة تتعلّق فوق خط عرض باريس سيكون مركزها
        على محور الأرض، شمال المركز. واستنتج: في أي
        مستوٍ يجب أن يقع \omterm{def:g12:satellites-kepler:geo}{مدار ثابت بالنسبة إلى الأرض}؟
  \item اكتب \omterm{thm:g12:newtons-laws:second}{قانون نيوتن الثاني} في \omterm{def:g12:kinematics-2d:frenet}{معلم فريني} من أجل \omterm{def:g10:universal-gravitation:satellite}{مدار}
        دائري نصف قطره $r$ واشتقّ $v = \sqrt{GM/r}$.
  \item استنتج $T = 2\pi\sqrt{r^3/(GM)}$ وأعِد كتابته على صورة قانون كبلر
        الثالث $T^2/r^3 = 4\pi^2/(GM)$.
\end{enumerate}

\medskip
\noindent\textbf{الجزء II --- الدائرة الوحيدة التي تفي.}
\begin{enumerate}[resume]
  \item حُلّ القانون الثالث لإيجاد $r$ واحسب نصف قطر \omterm{def:g12:satellites-kepler:geo}{المدار
        الثابت بالنسبة إلى الأرض}.
  \item استنتج الارتفاع فوق الأرض.
  \item احسب \omterm{def:g12:kinematics-2d:velocity}{السرعة} المدارية من $v = 2\pi r/T$.
  \item وتحقّق منها بالعلاقة $v = \sqrt{GM/r}$.
  \item وكتلة \omterm{def:g10:universal-gravitation:satellite}{تابع} الشركة \qty{4.0e3}{kg}؛ \omterm{def:g10:universal-gravitation:satellite}{وتابع} منافس
        كتلته ضعف ذلك. قارن نصفَي قطر مداريهما الثابتين بالنسبة إلى الأرض،
        وقل ما الذي تغيّره الكتلة الزائدة عند الإطلاق.
\end{enumerate}

\medskip
\noindent\textbf{الجزء III --- المحطة تحته.}
\begin{enumerate}[resume]
  \item تحلّق محطة الفضاء على $r = \qty{6.77e6}{m}$: احسب
        سرعتها.
  \item احسب دورها، بالثواني وبالدقائق.
  \item وكم مدارًا تكمل في \qty{24}{h} --- وكم شروقًا تقريبًا
        يشاهد روّاد الفضاء في اليوم؟
  \item احسب نسبة نصفَي القطر $r_{\text{الثابت}}/r_{\text{المحطة}}$،
        وتحقّق من أن نسبة الدورين هي ذلك العدد مرفوعًا إلى
        \omterm{def:g10:orders-of-magnitude:scinot}{الأس} $3/2$، كما يقتضي قانون كبلر الثالث.
  \item وأيّ التابعين أسرع، وأيّهما كلّف \omterm{def:g10:energy-conservation:energy}{طاقة} أكبر لكل
        \omterm{def:g10:orders-of-magnitude:unit}{كيلوغرام} لوضعه؟ ووفّق بينهما في جملة واحدة.
\end{enumerate}

\medskip
\noindent\textbf{الجزء IV --- \omterm{def:g10:universal-gravitation:weight}{وزن} المشتري على الهامش.}
\begin{enumerate}[resume]
  \item انطلاقًا من قانون كبلر الثالث مطبَّقًا على آيو، عبّر عن كتلة
        المشتري $M_J$ بدلالة $G$ و$r$ و$T$.
  \item احسب $M_J$.
  \item وكم كتلة أرضية هي؟ وأيّ جزء تقريبًا من
        كتلة الشمس؟
  \item ومرحّل ``ثابت بالنسبة إلى المشتري'' يجب أن يتعلّق فوق نقطة من
        سحب المشتري: احسب نصف قطر \omterm{def:g10:universal-gravitation:satellite}{مداره}؛ وقارنه
        \omterm{def:g10:universal-gravitation:satellite}{بمدار} آيو.
  \item ارفع تقريرًا إلى مجلس الإدارة، في جملتين: الارتفاع الذي يجب أن تركن فيه
        الشركة مرحّلها، وكتلة المشتري التي
        سلّمها القانون نفسه مجانًا.
\end{enumerate}
\end{problem}
